\documentclass[11pt,a4paper]{article}
\usepackage[T1]{fontenc}
\usepackage{lmodern,amsmath,amssymb,amsthm,mathtools,booktabs,array,microtype}
\usepackage[margin=27mm,headheight=14pt]{geometry}
\usepackage[authoryear,round]{natbib}
\usepackage{fancyhdr}
\usepackage[hidelinks]{hyperref}
\usepackage{xurl}
\pagestyle{fancy}\fancyhf{}\fancyhead[L]{\small Centroidal equilibrium creation under abrasion}\fancyhead[R]{\small Candidate v1.1}\fancyfoot[C]{\thepage}
\newtheorem{theorem}{Theorem}[section]
\newtheorem{lemma}[theorem]{Lemma}
\newtheorem{proposition}[theorem]{Proposition}
\newtheorem{corollary}[theorem]{Corollary}
\theoremstyle{remark}\newtheorem{remark}[theorem]{Remark}
\newcommand{\R}{\mathbb R}\newcommand{\E}{\mathbb E}\newcommand{\Prob}{\mathbb P}
\DeclareMathOperator{\tr}{tr}\DeclareMathOperator{\sgn}{sgn}\DeclareMathOperator{\Cov}{Cov}
\title{Creation of centroidal equilibria under curvature abrasion,\ and a moment obstruction to stochastic monotonicity}
\author{Anonymous}
\date{5 October 2026\quad\textnormal{Prospective Evidence Press companion, version 1.1}}
\begin{document}
\maketitle
\begin{center}\small Unrefereed candidate. DOI: \url{https://doi.org/10.5281/zenodo.23167004}\end{center}
\begin{abstract}
An explicit polynomial support-function family defines smooth, uniformly strictly convex, centrally symmetric bodies whose number of centroidal static equilibria increases from 22 to 26 under inward mean-curvature flow. The proof is analytic: exact global critical-point enumeration is combined with a forward-time saddle-node unfolding, not a discretization of a geometric PDE. The same family works for every nontrivial nonnegative combination of mean and Gaussian curvature speeds, with an arbitrary nonnegative constant speed added. Endpoint count increase persists on an open set of asymmetric initial bodies and yields counterexamples to distribution-free expected-count monotonicity. Separately, a four-outcome probability law satisfies the zero-mean and zero-covariance assumptions in Domokos's stochastic argument but reverses its asserted drift. We give the exact probability condition for nonpositive drift and sufficient conditional-symmetry hypotheses. The published centroid conjecture concerns an expected count, not deterministic monotonicity; the distinction and the limits of both counterexamples are explicit throughout.
\end{abstract}
\noindent\textbf{Keywords.} Geometric abrasion; mean-curvature flow; support function; static equilibrium; saddle-node bifurcation; covariance; stochastic drift.

\section{The two statements must be separated}
The supplied geology problem bundle asks whether the number of static equilibria, measured from a convex body's centroid, decreases monotonically under curvature-driven abrasion. Taken as a universal deterministic statement, the answer is negative. We prove this using the exact body family below.

The actual Conjecture 1 of \citet{D15}, however, is formulated for the expected value \(\bar N(t)\). The same paper already acknowledges that deterministic spatial critical-point counts can increase on surfaces, discusses a heat-equation birth example, and studies nonmonotonic centroidal counts in an illustrative planar model. We do not claim that the general possibility of creation is new. Our deterministic contribution is an explicit globally convex three-dimensional centroidal construction, with all equilibria enumerated and a genuine curvature-flow birth certified analytically.

The stochastic result in \citet[Assumption 2 and equations (49)--(53)]{D15} replaces third derivatives by random variables with zero means and covariance. Its claimed sign inference does not follow from those moments. Section~\ref{sec:probability} gives an exact counterexample and replacement conditions. This invalidates the moment-based drift conclusion, not every possible stochastic model of abrasion. In particular, an alternative sampling hypothesis called Assumption 2-A in that paper is logically distinct and is not refuted by our four-outcome law.

The general local creation mechanism under Gaussian smoothing already belongs to the scale-space literature, notably \citet{Damon95}. That result is background, not the theorem claimed here: the present construction controls a closed strictly convex body's centroid and enumerates all equilibria under the actual nonlinear curvature equation. The separate probability result repairs a moment-to-sign inference; it does not supply a physical sampling law.

\begin{theorem}[Explicit centroidal creation]\label{thm:main}
Put \(\varepsilon=1/100\). For \(u=(x,y,z)\in S^2\), define
\begin{equation}\label{eq:support}
 s_\mu(u)=1+\varepsilon f_\mu(u),\qquad
 f_\mu(u)=\frac{y^2}{2}+\frac{x^3z}{6}-2xy^2z+\mu xz.
\end{equation}
For \(|\mu|\le1\), this is the support function of a smooth strictly convex body with every principal curvature radius at least \(9/10\). The homogeneous body's centroid is the origin. For all sufficiently small \(\mu>0\), inward mean-curvature flow starting from this body undergoes two antipodal saddle-node births at a positive time \(t_*(\mu)\). In a neighborhood of this time, away from the bifurcation instant,
\[
 N(t)=22\quad(t<t_*),\qquad N(t)=26\quad(t>t_*),
 \qquad t_*(\mu)=\frac{10201}{29799}\mu+O(\mu^2).
\]
The centroid stays fixed throughout the flow. Initially and immediately before the births there are 6 stable equilibria, 6 unstable equilibria and 10 saddles; afterwards there are 8 stable equilibria, 6 unstable equilibria and 12 saddles.
\end{theorem}

The assertion ``sufficiently small'' is an analytic quantifier proved below, not a numerical guess at a parameter threshold. The package does not certify a particular numerical upper bound on \(\mu\) or integrate the flow numerically.

\section{Support geometry and uniform convexity}
For a smooth function \(s:S^2\to\R\), write
\[
 Q_s=\nabla^2_{S^2}s+sI.
\]
When \(s>0\) and \(Q_s\) is positive definite, the homogeneous extension \(\widetilde s(v)=|v|s(v/|v|)\) is convex: its tangential Hessian is \(Q_s/|v|\), its radial Hessian vanishes, and positivity gives convexity also on lines through the origin. It is therefore a support function. Its boundary in outward-normal coordinates is
\begin{equation}\label{eq:normalparam}
 X(u)=s(u)u+\nabla_{S^2}s(u),\qquad dX=Q_s,
\end{equation}
and the eigenvalues of \(Q_s\) are the principal curvature radii.

\begin{lemma}\label{lem:convex}
For the functions \eqref{eq:support}, \(s_\mu>0\) and \(Q_{s_\mu}\succeq(9/10)I\) whenever \(|\mu|\le1\).
\end{lemma}
\begin{proof}
For the restriction of an ambient polynomial \(f\) to the unit sphere,
\[
 Q_f=(D^2 f)|_{u^\perp}+(f-u\cdot\nabla f)I.
\]
Set \(f_2=y^2/2\), \(g=x^3z/6-2xy^2z\) and \(h=xz\). The ambient Hessian of \(g\) is
\[
 D^2g=\begin{pmatrix}
 xz&-4yz&x^2/2-2y^2\\
 -4yz&-4xz&-4xy\\
 x^2/2-2y^2&-4xy&0
 \end{pmatrix}.
\]
On \(S^2\), its maximum absolute row sum is at most 6, so its operator norm is at most 6. Also \(\|D^2f_2\|\le1\) and \(\|D^2h\|=1\). The bounds \(|f_2|\le1/2\), \(|h|\le1/2\), and
\[
 |g|\le\frac1{12}+\frac14=\frac13
\]
follow from \(|xz|\le1/2\) and \(2|xz|y^2\le(1-y^2)y^2\le1/4\). Homogeneity gives
\[
 f_\mu-u\cdot\nabla f_\mu=-f_2-3g-\mu h.
\]
Thus
\[
 \|Q_{f_\mu}\|\le 6+1+|\mu|+\frac32+\frac{|\mu|}{2}
 =\frac{17}{2}+\frac32|\mu|\le10.
\]
Consequently \(Q_{s_\mu}=I+\varepsilon Q_{f_\mu}\succeq(1-10\varepsilon)I=(9/10)I\). The bound \(|f_\mu|\le4/3\) also gives \(s_\mu>0\).
\end{proof}

Every term of \(f_\mu\) is even under \(u\mapsto-u\). Hence the body is centrally symmetric and its homogeneous centroid is zero. It is also invariant under \(y\mapsto-y\). Isometry invariance and uniqueness of the flow preserve both symmetries.

\begin{lemma}[Equilibrium correspondence]\label{lem:equilibria}
For a strictly convex body containing the origin, centroidal static equilibria are exactly the critical points of its support function, when its centroid is the origin. Their Morse types agree with those of the radial distance from the origin.
\end{lemma}
\begin{proof}
A horizontal supporting plane through \(X(u)\) balances the body precisely when the normal through the support point passes through the centroid. By \eqref{eq:normalparam}, \(X(u)\) is parallel to \(u\) exactly when \(\nabla s=0\). For the radial-distance type, differentiate
\[
 |X|^2=s^2+|\nabla s|^2,
 \qquad \nabla(|X|^2)=2Q_s\nabla s.
\]
At a critical point the Hessian is \(2Q_s\nabla^2s\). The factors commute because \(Q_s=sI+\nabla^2s\), and \(Q_s\) is positive definite, so the Hessians have the same inertia. The Gauss parametrization and radial parametrization are diffeomorphic, preserving the critical-point classification.
\end{proof}

\section{A forward-time fold, not a backward PDE argument}\label{sec:fold}
Here mean curvature means the \emph{sum} $\kappa_1+\kappa_2$, not its average. The averaged convention doubles the displayed birth time. For inward mean-curvature flow, the support equation is
\begin{equation}\label{eq:mcf}
 \partial_t s=-\tr(Q_s^{-1}).
\end{equation}
Smooth strictly convex initial surfaces have a unique smooth solution for a positive short time \citep{H84}. The following records precisely the local dependence needed here.

\begin{lemma}[Common forward interval and dependence]\label{lem:dependence}
Fix $0<\alpha<1$ and one smooth strictly convex initial support function. For the mean-curvature equation, or any fixed speed in Theorem~\ref{thm:bloore}, some $C^{6,\alpha}(S^2)$ neighborhood of that function has a common existence interval $[0,T]$. Its solution map is continuous into $C([0,T];C^2(S^2))$. For a smooth finite-dimensional family of smooth data, the solution depends smoothly on the parameters and has the spatial and forward-time derivatives used below, including at $t=0$.
\end{lemma}
\begin{proof}
On a small neighborhood $Q$ remains uniformly positive definite and bounded. The principal linearized coefficient is $Q^{-2}$, or $bQ^{-2}+c(\det Q)^{-1}Q^{-1}$; it is uniformly positive for fixed $b,c\ge0$, $b+c>0$. The nonlinear operator is smooth there. Apply the closed-manifold local existence theorem and linear Schauder isomorphism of \citet[Theorems 1.1--1.2 and 2.3--2.4]{Huang}. On a fixed short cylinder, the map $s\mapsto(s_t-F(s),s(0))$ has invertible derivative in parabolic H\"older spaces. The Banach implicit-function theorem therefore supplies the same interval for nearby data and smooth parameter dependence. Spatial bootstrap gives the stated derivatives for smooth data; $C^{6,\alpha}$ control is more than needed for endpoint $C^2$ continuity. This is entirely a forward-time argument.
\end{proof}
The lemma applies near $s_0$ to the $\mu$ family and near the fixed positive-$\mu$ datum to the asymmetric perturbations in Corollary~\ref{cor:robust}.

At \(p=(0,0,1)\), use the chart \((x,y,\sqrt{1-x^2-y^2})\). At \(t=\mu=0\), the exact jets are
\begin{equation}\label{eq:jets}
 \begin{gathered}
 s=1,\quad s_x=s_y=s_{xx}=s_{xy}=0,\
 s_{yy}=\varepsilon,\qquad s_{xxx}=\varepsilon,
 \qquad s_{xyy}=-4\varepsilon.
 \end{gathered}
\end{equation}
The Christoffel symbols vanish at this chart origin. Because \(\nabla s=0\), the relevant third partial derivatives equal the covariant third derivatives there. Thus \(Q_s=\operatorname{diag}(1,1+\varepsilon)\) at the pole, and differentiating \eqref{eq:mcf} gives
\begin{equation}\label{eq:sxt}
 s_{xt}=\tr(Q_s^{-1}(\partial_xQ_s)Q_s^{-1})
 =\varepsilon\left(1-\frac4{(1+\varepsilon)^2}\right)
 =-\frac{29799}{1020100}<0.
\end{equation}
This has the opposite sign from the positive cubic derivative in \eqref{eq:jets}, the orientation for a birth.

To make the event occur at positive time from nondegenerate initial data, let \(s(x,y,t;\mu)\) be the forward solution and put
\[
 G(x,t,\mu)=s_x(x,0,t;\mu).
\]
Reflection symmetry gives \(s_y(x,0,t;\mu)=0\). At \((x,t,\mu)=(0,0,0)\),
\[
 \begin{gathered}
 G=G_x=0,\quad G_{xx}=\varepsilon,\quad
 G_\mu=\varepsilon,\quad G_{x\mu}=0,\\
 G_t=-\varepsilon k,\qquad
 k=\frac4{(1+\varepsilon)^2}-1=\frac{29799}{10201}>0.
 \end{gathered}
\]
The derivative of \((G,G_x)\) with respect to \((\mu,x)\) is \(\operatorname{diag}(\varepsilon,\varepsilon)\), hence invertible. The parameter-dependent implicit-function theorem, applied for forward \(t\ge0\), supplies the fold curve
\[
 \mu_*(t)=kt+O(t^2),\qquad x_*(t)=O(t).
\]
One may apply the theorem at each forward time in a common local neighborhood; no negative-time solution is required. Since \(\mu_*'(0)=k>0\), the curve can be inverted:
\[
 t_*(\mu)=\mu/k+O(\mu^2)>0\qquad(0<\mu\text{ sufficiently small}).
\]

The transverse derivative \(s_{yy}\) remains positive. By the implicit-function theorem the only nearby solution of \(s_y=0\) is \(y=0\). The function \(G\) has positive second \(x\)-derivative and negative \(t\)-derivative at its fold. Its local minimum crosses zero downwards: there are no local critical points immediately before the fold and two immediately afterwards, one minimum and one saddle. The same is true at the antipode by central symmetry. At \(t=0\) and small positive \(\mu\), the minimum of \(G\) is \(\varepsilon\mu+O(\mu^2)>0\), so the initial body has no critical point in either pole neighborhood. Section~\ref{sec:global} proves that exactly 22 other, Morse, points persist and that no other births interfere on the common short interval. This completes the fold part of Theorem~\ref{thm:main}.

\section{Exact enumeration of every initial critical point}\label{sec:global}
The critical points of \(s_0\) are those of \(f_0\). Parameterize away from \(y=\pm1\) by
\[
 y=r,\qquad x=\sqrt{1-r^2}\sin\theta,\qquad
 z=\sqrt{1-r^2}\cos\theta.
\]
Let \(q=r^2\), \(t=\tan\theta\), and
\[
 A(t)=\frac{t^3}{6(1+t^2)^2},\qquad B(t)=\frac{t}{1+t^2}.
\]
Then
\begin{equation}\label{eq:reduced}
 f_0=F(t,q)=A+q(1/2-2A-2B)+q^2(A+2B).
\end{equation}
We handle the coordinate exceptions explicitly.

On \(y=0\),
\[
 \frac{d}{d\theta}\left(\frac{\sin^3\theta\cos\theta}{6}\right)
 =\frac{\sin^2\theta}{6}(3\cos^2\theta-\sin^2\theta).
\]
There are the two degenerate points \((0,0,\pm1)\) and four Morse points with \(\tan\theta=\pm\sqrt3\). At \(t=\pm\sqrt3\), the transverse second derivative is \(1\mp9\sqrt3/8\), and the angular second derivative has the same sign. These four points comprise two minima and two maxima.

At \((0,\pm1,0)\), the Hessian in local \((x,z)\) coordinates is
\[
 \begin{pmatrix}-1&-2\\-2&-1\end{pmatrix},
\]
so these are two saddles. At \(z=0\) with \(0<|y|<1\), the derivative \(F_q\) is \(1/2\); there are no missing critical points on that chart boundary.

For the remaining points, \(0<q<1\) and \(t\) is finite. The equation \(F_t=0\) yields
\begin{equation}\label{eq:qstar}
 q=q_*(t)=\frac{t^2(t^2-3)}{D(t)},\qquad D(t)=13t^4-3t^2-12.
\end{equation}
If \(D(t)=0\), the numerator that must vanish in \(F_t\) cannot vanish, so no solution is lost by division. Substituting \eqref{eq:qstar} into \(F_q=0\) gives
\begin{equation}\label{eq:quartic}
 p(t)=13t^4-52t^3-3t^2+48t-12=0,
 \qquad F_q(t,q_*(t))=\frac{p(t)}{2D(t)}.
\end{equation}
The discriminant is \(25315198272\), and \(p\) is relatively prime to \(D\). Exact Sturm counts isolate its four simple real roots as follows; each root corresponds to four sphere points, from the two signs of \(y\) and the two antipodal choices of \((x,z)\).
\begin{center}
\begin{tabular}{llll}
\toprule
Rational root interval & Roots & Sphere points & Type\\
\midrule
\((-99/100,-98/100)\)&1&4&Maximum\\
\((27/100,28/100)\)&1&4&Saddle\\
\((88/100,89/100)\)&1&4&Minimum\\
\((382/100,383/100)\)&1&4&Saddle\\
\bottomrule
\end{tabular}
\end{center}
This table concerns the limiting initial surface at $\mu=0$, not a numerically sampled positive-time trajectory.
In each interval \(0<q_*(t)<1\). The exact Hessian determinant in \((t,r)\) coordinates at a critical point is
\begin{equation}\label{eq:hessdet}
 -\frac{16t^2(t-1)(t+1)(t^2-3)
 (169t^6-429t^4+504t^2-144)}
 {(1+t^2)^2(13t^4-3t^2-12)^3}.
\end{equation}
The transverse Hessian sign is the sign of
\(F_{qq}=t(13t^2+12)/(3(1+t^2)^2)\), hence the sign of \(t\). Sturm counts certify that no numerator or denominator affecting the classification vanishes in any listed interval. Evaluation at a rational midpoint then certifies all the displayed signs.

Adding the exceptional charts gives exactly 6 minima, 6 maxima and 10 saddles, together with the two degenerate pole points. The Morse counts satisfy \(6+6-10=2\), as a consistency check. Compactness bounds the gradient away from zero outside fixed small neighborhoods of these 24 points. The 22 Morse points persist under sufficiently small changes of \(\mu\) and sufficiently short forward evolution. The two remaining neighborhoods are completely described by Section~\ref{sec:fold}. The counts before and after the births are therefore exactly 22 and 26, proving Theorem~\ref{thm:main}.

\section{Extensions to abrasion speeds and nonsymmetric ensembles}
\begin{theorem}[Nonnegative Bloore-type combinations]\label{thm:bloore}
Let
\[
 v=a+b(\kappa_1+\kappa_2)+c\kappa_1\kappa_2,
 \qquad a,b,c\ge0,\quad b+c>0,
\]
be the inward normal speed. The family \eqref{eq:support}, for all sufficiently small positive \(\mu\), again has a centroidal equilibrium increase from 22 to 26 at positive time. The allowed \(\mu\) interval may depend on \(a,b,c\).
\end{theorem}
\begin{proof}
The support equation is
\[
 s_t=-a-b\tr(Q_s^{-1})-c(\det Q_s)^{-1}.
\]
At the pole the null-direction gradient derivative is
\begin{equation}\label{eq:bloorejet}
 s_{xt}=\varepsilon\left[
 b\left(1-\frac4{(1+\varepsilon)^2}\right)
 +c\left(\frac1{1+\varepsilon}-\frac4{(1+\varepsilon)^2}\right)
 \right]<0.
\end{equation}
Both brackets are negative; the Gaussian contribution, including \(\varepsilon\), is \(-299/10201\). The constant speed has zero spatial derivative. The equation is strictly parabolic near the initial data, because \(v_{\kappa_i}=b+c\kappa_j>0\). Equivalently, the derivative of its support operator in the curvature-radius matrix is a positive combination of \(Q^{-2}\) and \((\det Q)^{-1}Q^{-1}\). Local existence, uniqueness and parameter dependence follow from the same uniformly parabolic theory. Every remaining part of the preceding argument applies unchanged, with \(k=-s_{xt}/\varepsilon>0\).
\end{proof}

This includes mean-curvature and Gaussian-curvature flow and the spherical-abrader speed \((1+R\kappa_1)(1+R\kappa_2)\) with \(R>0\), discussed in \citet{D15}. It does not include pure constant-speed flow \(b=c=0\), for which the strict-parabolic and local-birth arguments used here do not apply.

\begin{corollary}[Open-set and ensemble robustness]\label{cor:robust}
For each fixed flow in Theorem~\ref{thm:bloore}, there are times \(0<\tau_-<\tau_+\) and an open set of smooth strictly convex initial bodies, including asymmetric bodies, whose centroidal equilibrium counts satisfy
\[
 N(\tau_-)=22,\qquad N(\tau_+)=26.
\]
Any probability measure supported on this open set has an expected-count increase between these times. A nondegenerate example is the uniform law on a sufficiently small positive parameter interval in \eqref{eq:support}.
\end{corollary}
\begin{proof}
Fix a sufficiently small \(\mu>0\), and choose \(\tau_-,\tau_+\) strictly before and after its birth time, while all equilibria at the two endpoints are Morse. Solutions depend continuously on their smooth initial support functions, on a common time interval and in the \(C^2\) norm at these endpoints. The homogeneous centroid also depends continuously on the body; recentering replaces \(s(u)\) by \(s(u)-G\cdot u\). Thus the centered support functions at the two endpoints remain \(C^2\)-close for sufficiently small perturbations of the initial body. Morse critical-point counts on the compact sphere are stable under such perturbations: their Hessians stay invertible locally and the gradient stays bounded away from zero elsewhere. The counts are therefore fixed at 22 and 26 on an open neighborhood in a sufficiently high smooth norm. This neighborhood contains asymmetric bodies, for which the centroid need not stay fixed. Integrating the two constant endpoint counts against any supported probability measure proves the expectation claim. Continuity of the parameter family gives a nonzero-length interval of positive \(\mu\) values inside this neighborhood.
\end{proof}

This disproves a \emph{distribution-free} expectation claim. It is not a claim that the ensemble in Corollary~\ref{cor:robust} satisfies the random-jet moment assumptions of \citet{D15}, or that it represents naturally sampled rocks.

\section{An exact obstruction to the stochastic moment argument}\label{sec:probability}
In the notation of \citet[equations (49)--(53)]{D15}, put
\[
 A=\bar r_{yyy},\quad B=\bar r_{xxy},\quad
 \alpha=v_\kappa>0,\quad\beta=v_\lambda>0.
\]
The annihilation indicator and jump size are
\begin{equation}\label{eq:sign}
 \Omega=\sgn(\alpha A^2+\beta AB),\qquad \Delta N=-2\Omega.
\end{equation}
Assumption 2 is \(\E A=\E B=\Cov(A,B)=0\). This gives a positive expectation for the expression inside the sign, provided \(A\) has nonzero variance, but does not determine the expected sign.

\begin{proposition}[Four-outcome counterexample]\label{prop:law}
Even when \(\alpha=\beta=1\), both marginal laws are symmetric, all moments are finite, and the three moment conditions in Assumption 2 hold, one can have
\[
 \E(A^2+AB)=1>0,\qquad \E\Omega=-\frac45,
 \qquad\E\Delta N=\frac85>0.
\]
\end{proposition}
\begin{proof}
Use this probability law:
\begin{center}
\begin{tabular}{rrrr}
\toprule
Probability & \(A\)&\(B\)&\(A^2+AB\)\\
\midrule
\(9/20\)&1&$-2$&$-1$\\
\(9/20\)&$-1$&2&$-1$\\
\(1/20\)&1&18&19\\
\(1/20\)&$-1$&$-18$&19\\
\bottomrule
\end{tabular}
\end{center}
The paired outcomes give \(\E A=\E B=0\). Also
\[
 \E AB=\frac9{10}(-2)+\frac1{10}(18)=0,
 \quad \E A^2=1,\quad \E B^2=36.
\]
Yet the inside expression is negative with probability \(9/10\), positive with probability \(1/10\), and never zero. Thus \(\E\Omega=1/10-9/10=-4/5\), and \eqref{eq:sign} gives the positive mean jump.
\end{proof}

The step from equation (50) to (51) in \citet{D15} is therefore false under its stated moment hypothesis, and the resulting embedded jump-chain drift is not determined by those moments. The example also survives the addition of an independent uniform noise \(\eta\in[-d,d]\) for any \(0<d<1\): neither of the values \(-1\) and 19 changes sign. Thus adding small independent centered noise to this moment model does not repair the inference. This observation does not identify \(AB\) with the alternative determinant-sampling variable in Assumption 2-A.

\begin{proposition}[The sign can vary across its full range]\label{prop:sharp}
Under the same zero-mean and zero-covariance constraints, with \(\E A^2=1\), \(\E\sgn(A^2+AB)\) can take every value in \((-1,1)\).
\end{proposition}
\begin{proof}
Let \(A\) be a symmetric Rademacher variable, independent of \(R\). For any \(q\in(0,1)\), take
\[
 \Prob(R=-2)=q,\qquad
 \Prob\left(R=\frac{2q}{1-q}\right)=1-q,
 \qquad B=RA.
\]
Then \(\E R=0\), hence \(\E A=\E B=\E AB=0\), while \(A^2+AB=1+R\). Its expected sign is \(1-2q\) and its expectation is 1. Every member of this family has finite variance; no common variance bound is asserted as \(q\to1\).
\end{proof}

\subsection{A correct probability criterion and sufficient hypotheses}
\begin{theorem}[Exact drift criterion]\label{thm:criterion}
Suppose \(A\ne0\) almost surely and \(\Prob(\alpha A^2+\beta AB=0)=0\), with \(\alpha,\beta>0\). Then
\begin{equation}\label{eq:criterion}
 \E\Delta N=-2+4\Prob\{A(\alpha A+\beta B)<0\}.
\end{equation}
Consequently nonpositive mean jump is equivalent to the probability on the right being at most \(1/2\). A sufficient condition is symmetry of the conditional law of \(B\) given \(A\), around zero. Independence of \(A,B\), together with symmetry of \(B\), is a stronger sufficient condition.
\end{theorem}
\begin{proof}
There are only the two signs, so \(\E\Omega=1-2\Prob(\Omega=-1)\), proving \eqref{eq:criterion}. Conditional on \(A=a>0\), a negative sign requires \(B<-\alpha a/\beta\). Conditional on \(A=a<0\), it requires \(B>\alpha|a|/\beta\). Under conditional symmetry each probability is at most \(1/2\). Integrating proves the claim. Strictly negative drift needs additional conditional mass inside the corresponding threshold intervals; symmetry alone can yield zero drift.
\end{proof}

For example, independent symmetric variables \(A=\pm1\) and \(B=\pm2\), with \(\alpha=\beta=1\), have zero mean drift rather than the 75\% annihilation probability quoted under symmetry in equation (52) of \citet{D15}. Identical distribution is a useful additional sufficient hypothesis: if \(A,B\) are independent, identically distributed, continuous and symmetric, the event \(AB<0\) has probability \(1/2\), independently of their magnitudes, and exchangeability gives \(\Prob(|B|>|A|)=1/2\). Thus the creation probability is \(1/4\), and the annihilation probability is \(3/4\), when \(\alpha=\beta\).

More generally, independent centered Gaussians with positive standard deviations \(\sigma_A,\sigma_B\) give the explicit formula
\[
 \E\Omega=\frac2\pi\arctan\left(\frac{\alpha\sigma_A}{\beta\sigma_B}\right),
 \qquad
 \E\Delta N=-\frac4\pi\arctan\left(\frac{\alpha\sigma_A}{\beta\sigma_B}\right)<0.
\]
Indeed, after standardization the joint normal vector has uniform polar angle, so its coordinate ratio is standard Cauchy. The sign is negative precisely when \(B/A<-\alpha/\beta\). Integrating the Cauchy density gives the formula. This also shows why the individual variances can matter, despite their absence from the moment condition.

A continuous-time expected-count theorem additionally needs conditional drift conditions at each state and a well-defined nonexplosive event-rate/holding-time model. Equation~\eqref{eq:criterion} resolves the local embedded-jump question; it does not prescribe such a model or an empirical rock ensemble.

\section{Verification, dependencies and scope}
The script \path{code/verify_abrasion.py} uses exact SymPy algebra and rational probability arithmetic. It verifies the support-function symmetry, the displayed jets, both curvature-flow signs, the rational convexity majorant, the critical-point reduction, Sturm root counts, all Hessian classifications, the final 22/26 counts and the four-outcome moments. The Gaussian and conditional-symmetry statements above are proved directly, not by simulation.

The continuum conclusion also uses smooth short-time existence, uniqueness and parameter dependence for uniformly parabolic curvature flows, and the implicit-function and Morse-persistence arguments given here. These analytic steps are not replaced by finite sampling and are not formalized in a proof assistant. The supplied verification report is not an independent human referee report.

The resolved claims are: deterministic centroidal monotonicity is false even for a smooth strictly convex centrally symmetric body under mean-curvature flow; the failure extends to the stated Bloore-type speeds and an open set of asymmetric bodies; a distribution-free expected-count version is false; and the stated moment assumptions do not imply the published drift sign. A particular fully specified stochastic abrasion ensemble may still have decreasing expected counts under stronger assumptions. No assertion about observed frequencies in natural rocks follows from these mathematical counterexamples.

\begin{thebibliography}{9}
\bibitem[Damon(1995)]{Damon95}
Damon, J. (1995). Local Morse theory for solutions to the heat equation and Gaussian blurring. \emph{Journal of Differential Equations, 115}(2), 368--401. \url{https://doi.org/10.1006/jdeq.1995.1019}.

\bibitem[Huang(2026)]{Huang}
Huang, H. (2026). The Cauchy problem for fully nonlinear parabolic systems on manifolds. arXiv:1506.05030v8, 25 August 2026; first posted 2015. \url{https://arxiv.org/abs/1506.05030v8}.

\bibitem[Domokos(2015)]{D15}
Domokos, G. (2015). Monotonicity of spatial critical points evolving under curvature-driven flows. \emph{Journal of Nonlinear Science, 25}, 247--275. Published online 2014. \url{https://doi.org/10.1007/s00332-014-9228-3}. Author version, v4: \url{https://arxiv.org/abs/1308.4779}.

\bibitem[Huisken(1984)]{H84}
Huisken, G. (1984). Flow by mean curvature of convex surfaces into spheres. \emph{Journal of Differential Geometry, 20}(1), 237--266. \url{https://doi.org/10.4310/jdg/1214438998}. Related author lecture: \url{https://maths.anu.edu.au/files/CMAProcVol8-Huisken_1.pdf}.
\end{thebibliography}
\end{document}
